% !TEX root = ../../../main.tex
% Block card: catalogue.manual_labels -- "Manual labels" (Part I, the blocks chapter).
\cardfield{What it does}
Gives each window the label a person gave that stretch of signal (interesting or not
interesting), so a classifier can learn from human judgement.

\cardfield{When to use it}
In the human-label arm of research question 1, in place of \blockref{catalogue.cluster}: the
classifier after it trains on these labels exactly as it would on clusters. It reads the labels
you made in Review and Explore and the labels imported from the catalogue.

\begin{plainterms}[When does a window take a label?]
A label covers a stretch of time. A window takes it only if the window holds the whole labelled
stretch, like a box that only counts as holding a parcel when the parcel is entirely inside it.
If a window holds two labelled stretches that disagree, it is set aside as conflicting. A window
lying inside a long ``not interesting'' stretch is not interesting, since nothing is anywhere in
it; inside a long ``interesting'' stretch it stays unlabelled, since the event could be anywhere.
Time nobody labelled is unlabelled, never ``not interesting''.
\end{plainterms}

\cardfield{What goes in, what comes out}
In go windows. Out comes, per window, 1 (interesting), 0 (not interesting) or $-1$ (left out:
unlabelled, conflicting, artifact, or dropped to avoid overlap), with counts of each.

\cardfield{Settings that matter}
\emph{non\_overlapping} (on by default) keeps no two overlapping windows, as a training set
needs: labelled windows are kept first, then unlabelled ones fill the gaps.

\cardfield{Watch out for}
Only windows on the labels' own grid can take labels from the 10-minute set: a window off the
grid only ever partly covers a label, and the block says so. The block only reads: human labels
and machine detections live in separate tables, and nothing here writes either.
